\documentclass[10pt,a4paper]{amsart}
\usepackage{amsmath,amssymb,amsthm,graphicx}
% Retained preamble for the Stacks Project chapter {CH} (shared project preamble).
% The chapter file pinned-original-source.tex only carries \input{{preamble}} on its first line,
% so no preamble text lives inside the chapter itself.  The Stacks Project uses one single shared
% preamble for every chapter: the snapshot stages it three times (sources/stacks-chow/preamble.tex,
% sources/stacks-etale-cohomology/preamble.tex and focused-seeds/AGP/sources/stacks-snapshot-20261011/
% preamble.tex) and those copies are identical in size.  The block below is copied mechanically,
% line for line, from sources/stacks-chow/preamble.tex lines 14-18 and 154-217: the Xy-pic setup and
% the native theorem declarations together with the mathematical macros.  Nothing was added, removed
% or redefined; the wrapper class and the amsmath/amssymb/amsthm/graphicx base packages are supplied
% by the assembler.
% For commutative diagrams we use Xy-pic
% Chapter: stacks-cohomology.tex -- Cohomology of Sheaves.
\usepackage[all]{xy}

% We use 2cell for 2-commutative diagrams.
\xyoption{2cell}

% Theorem environments.
%
\theoremstyle{plain}
\newtheorem{theorem}[subsection]{Theorem}
\newtheorem{proposition}[subsection]{Proposition}
\newtheorem{lemma}[subsection]{Lemma}

\theoremstyle{definition}
\newtheorem{definition}[subsection]{Definition}
\newtheorem{example}[subsection]{Example}
\newtheorem{exercise}[subsection]{Exercise}
\newtheorem{situation}[subsection]{Situation}

\theoremstyle{remark}
\newtheorem{remark}[subsection]{Remark}
\newtheorem{remarks}[subsection]{Remarks}

\numberwithin{equation}{subsection}

% Macros
%
\def\lim{\mathop{\mathrm{lim}}\nolimits}
\def\colim{\mathop{\mathrm{colim}}\nolimits}
\def\Spec{\mathop{\mathrm{Spec}}}
\def\Hom{\mathop{\mathrm{Hom}}\nolimits}
\def\Ext{\mathop{\mathrm{Ext}}\nolimits}
\def\SheafHom{\mathop{\mathcal{H}\!\mathit{om}}\nolimits}
\def\SheafExt{\mathop{\mathcal{E}\!\mathit{xt}}\nolimits}
\def\Sch{\mathit{Sch}}
\def\Mor{\mathop{\mathrm{Mor}}\nolimits}
\def\Ob{\mathop{\mathrm{Ob}}\nolimits}
\def\Sh{\mathop{\mathit{Sh}}\nolimits}
\def\NL{\mathop{N\!L}\nolimits}
\def\CH{\mathop{\mathrm{CH}}\nolimits}
\def\proetale{{pro\text{-}\acute{e}tale}}
\def\etale{{\acute{e}tale}}
\def\QCoh{\mathit{QCoh}}
\def\Ker{\mathop{\mathrm{Ker}}}
\def\Im{\mathop{\mathrm{Im}}}
\def\Coker{\mathop{\mathrm{Coker}}}
\def\Coim{\mathop{\mathrm{Coim}}}

% Boxtimes
%
\DeclareMathSymbol{\boxtimes}{\mathbin}{AMSa}{"02}

%
% Macros for moduli stacks/spaces
%
\def\QCohstack{\mathcal{QC}\!\mathit{oh}}
\def\Cohstack{\mathcal{C}\!\mathit{oh}}
\def\Spacesstack{\mathcal{S}\!\mathit{paces}}
\def\Quotfunctor{\mathrm{Quot}}
\def\Hilbfunctor{\mathrm{Hilb}}
\def\Curvesstack{\mathcal{C}\!\mathit{urves}}
\def\Polarizedstack{\mathcal{P}\!\mathit{olarized}}
\def\Complexesstack{\mathcal{C}\!\mathit{omplexes}}
% \Pic is the operator that assigns to X its picard group, usage \Pic(X)
% \Picardstack_{X/B} denotes the Picard stack of X over B
% \Picardfunctor_{X/B} denotes the Picard functor of X over B
\def\Pic{\mathop{\mathrm{Pic}}\nolimits}
\def\Picardstack{\mathcal{P}\!\mathit{ic}}
\def\Picardfunctor{\mathrm{Pic}}
\def\Deformationcategory{\mathcal{D}\!\mathit{ef}}

% --- appended 2026-10-12 (Chief): the project own definitions of the reference/slogan/history
% display environments, copied mechanically from the shared Stacks preamble lines 8-12
% (\usepackage{verbatim} plus the three \newenvironment lines, which the chapter preamble keeps
% verbatim from the project preamble). Nothing added, removed or redefined: these environments
% carry no mathematical content and the project itself typesets them as comments. Without these
% three lines a counted statement whose own text uses them fails with Environment slogan undefined.
\usepackage{verbatim}
\newenvironment{reference}{\comment}{\endcomment}
\newenvironment{slogan}{\comment}{\endcomment}
\newenvironment{history}{\comment}{\endcomment}
\usepackage{cleveref}
\crefname{originalreference}{original reference}{original references}
\Crefname{originalreference}{Original reference}{Original references}
\makeatletter
\newlabel{categories-definition-dual}{{0FFP}{source}{}{sourcepointer-1}{}}
\newlabel{categories-definition-dual@cref}{{[originalreference][0][]0FFP}{[1][0][]source}{}{}{}}
\makeatother
\crefname{originalreference}{original reference}{original references}
\Crefname{originalreference}{Original reference}{Original references}
\makeatletter
\newlabel{example-dual-derived}{{0FPC}{source}{}{sourcepointer-2}{}}
\newlabel{example-dual-derived@cref}{{[originalreference][0][]0FPC}{[1][0][]source}{}{}{}}
\makeatother
\crefname{originallemma}{original lemma}{original lemmas}
\Crefname{originallemma}{Original lemma}{Original lemmas}
\makeatletter
\newlabel{lemma-K-flat-resolution}{{06YF}{source}{}{sourcepointer-3}{}}
\newlabel{lemma-K-flat-resolution@cref}{{[originallemma][0][]06YF}{[1][0][]source}{}{}{}}
\makeatother
\crefname{originalreference}{original reference}{original references}
\Crefname{originalreference}{Original reference}{Original references}
\makeatletter
\newlabel{modules-lemma-module-quotient-flat}{{05NI}{source}{}{sourcepointer-4}{}}
\newlabel{modules-lemma-module-quotient-flat@cref}{{[originalreference][0][]05NI}{[1][0][]source}{}{}{}}
\makeatother
\crefname{originalreference}{original reference}{original references}
\Crefname{originalreference}{Original reference}{Original references}
\makeatletter
\newlabel{derived-lemma-subcategory-left-resolution}{{05T7}{source}{}{sourcepointer-5}{}}
\newlabel{derived-lemma-subcategory-left-resolution@cref}{{[originalreference][0][]05T7}{[1][0][]source}{}{}{}}
\makeatother
\crefname{originallemma}{original lemma}{original lemmas}
\Crefname{originallemma}{Original lemma}{Original lemmas}
\makeatletter
\newlabel{lemma-perfect-precise}{{08CP}{source}{}{sourcepointer-6}{}}
\newlabel{lemma-perfect-precise@cref}{{[originallemma][0][]08CP}{[1][0][]source}{}{}{}}
\makeatother
\crefname{originalreference}{original reference}{original references}
\Crefname{originalreference}{Original reference}{Original references}
\makeatletter
\newlabel{derived-lemma-negative-exts}{{06XS}{source}{}{sourcepointer-7}{}}
\newlabel{derived-lemma-negative-exts@cref}{{[originalreference][0][]06XS}{[1][0][]source}{}{}{}}
\makeatother
\crefname{originallemma}{original lemma}{original lemmas}
\Crefname{originallemma}{Original lemma}{Original lemmas}
\makeatletter
\newlabel{lemma-local-actual}{{08C9}{source}{}{sourcepointer-8}{}}
\newlabel{lemma-local-actual@cref}{{[originallemma][0][]08C9}{[1][0][]source}{}{}{}}
\makeatother
\title{Left duals in the derived category are perfect and termwise dual}
\author{The Stacks Project Authors}
\date{stacks-project:cohomology.tex:local-snapshot-20261011; GNU FDL 1.2 or later}
\begin{document}
\maketitle
\noindent\textit{Editorial note.} All delivered mathematical text is the author's own, copied verbatim from the pinned Stacks Project chapter file cohomology.tex; no statement or proof step was rewritten, reordered or removed.\par\medskip
\noindent\textit{Editorial note.} The counted span cites the lemmas recorded in assembly.external\_references with canonical tags as navigation pointers.\par\medskip
\noindent\textit{Allowed standard prerequisites (editorial).}\begin{itemize}
\item K-flat resolutions. Cohomology of Sheaves, Lemma cohomology-lemma-K-flat-resolution. K-flat resolution of the left dual with flat terms.
\item Local representability of derived maps. Cohomology of Sheaves, Lemma cohomology-lemma-local-actual. sigma as an actual map of complexes, null-homotopy.
\item Precise perfection criterion. Cohomology of Sheaves, Lemma cohomology-lemma-perfect-precise. Concludes perfection from the descent.
\item Maps from complexes concentrated in negative degrees. Derived Categories, Lemma derived-lemma-negative-exts. L -> E\textasciicircum{}*[1] as a map of complexes.
\end{itemize}\medskip
\ifdefined\hypertarget\hypertarget{sourcepointer-1}{}\fi
\noindent\textit{Source reference.} Original reference 0FFP: sources/stacks-chow/tags.txt, source lines 19294--19294. Cross-chapter standard prerequisite cited by the author inside the counted span; it is used as a pointer and its own proof is not part of this item.\par\medskip
\ifdefined\hypertarget\hypertarget{sourcepointer-2}{}\fi
\noindent\textit{Source reference.} Original reference 0FPC: sources/stacks-chow/tags.txt, source lines 19597--19597. Cross-chapter standard prerequisite cited by the author inside the counted span; it is used as a pointer and its own proof is not part of this item.\par\medskip
\ifdefined\hypertarget\hypertarget{sourcepointer-3}{}\fi
\noindent\textit{Source reference.} Original lemma 06YF: sources/stacks-chow/tags.txt, source lines 8884--8884. Cross-chapter standard prerequisite cited by the author inside the counted span; it is used as a pointer and its own proof is not part of this item.\par\medskip
\ifdefined\hypertarget\hypertarget{sourcepointer-4}{}\fi
\noindent\textit{Source reference.} Original reference 05NI: sources/stacks-chow/tags.txt, source lines 7307--7307. Cross-chapter standard prerequisite cited by the author inside the counted span; it is used as a pointer and its own proof is not part of this item.\par\medskip
\ifdefined\hypertarget\hypertarget{sourcepointer-5}{}\fi
\noindent\textit{Source reference.} Original reference 05T7: sources/stacks-chow/tags.txt, source lines 7471--7471. Cross-chapter standard prerequisite cited by the author inside the counted span; it is used as a pointer and its own proof is not part of this item.\par\medskip
\ifdefined\hypertarget\hypertarget{sourcepointer-6}{}\fi
\noindent\textit{Source reference.} Original lemma 08CP: sources/stacks-chow/tags.txt, source lines 10609--10609. Cross-chapter standard prerequisite cited by the author inside the counted span; it is used as a pointer and its own proof is not part of this item.\par\medskip
\ifdefined\hypertarget\hypertarget{sourcepointer-7}{}\fi
\noindent\textit{Source reference.} Original reference 06XS: sources/stacks-chow/tags.txt, source lines 8861--8861. Cross-chapter standard prerequisite cited by the author inside the counted span; it is used as a pointer and its own proof is not part of this item.\par\medskip
\ifdefined\hypertarget\hypertarget{sourcepointer-8}{}\fi
\noindent\textit{Source reference.} Original lemma 08C9: sources/stacks-chow/tags.txt, source lines 10594--10594. Cross-chapter standard prerequisite cited by the author inside the counted span; it is used as a pointer and its own proof is not part of this item.\par\medskip
\begin{lemma}
\label{lemma-left-dual-derived}
Let $(X, \mathcal{O}_X)$ be a ringed space. Let $M$ be an object
of $D(\mathcal{O}_X)$. If $M$ has a left dual in the monoidal category
$D(\mathcal{O}_X)$ (Categories, Definition \ref{categories-definition-dual})
then $M$ is perfect and the left dual is as constructed in
Example \ref{example-dual-derived}.
\end{lemma}

\begin{proof}
Let $x \in X$. It suffices to find an open neighbourhood $U$ of $x$
such that $M$ restricts to a perfect complex over $U$. Hence during the
proof we can (finitely often) replace $X$ by an open neighbourhood of $x$.
Let $N, \eta, \epsilon$ be a left dual.

\medskip\noindent
We are going to use the following argument several times. Choose any
complex $\mathcal{M}^\bullet$
of $\mathcal{O}_X$-modules representing $M$. Choose a K-flat complex
$\mathcal{N}^\bullet$ representing $N$ whose terms are flat
$\mathcal{O}_X$-modules, see Lemma \ref{lemma-K-flat-resolution}.
Consider the map
$$
\eta : \mathcal{O}_X \to
\text{Tot}(\mathcal{M}^\bullet \otimes_{\mathcal{O}_X} \mathcal{N}^\bullet)
$$
After shrinking $X$ we can find an integer $N$ and for
$i = 1, \ldots, N$ integers $n_i \in \mathbf{Z}$ and sections
$f_i$ and $g_i$ of $\mathcal{M}^{n_i}$ and $\mathcal{N}^{-n_i}$
such that
$$
\eta(1) = \sum\nolimits_i f_i \otimes g_i
$$
Let $\mathcal{K}^\bullet \subset \mathcal{M}^\bullet$ be any subcomplex
of $\mathcal{O}_X$-modules containing the sections $f_i$
for $i = 1, \ldots, N$.
Since
$\text{Tot}(\mathcal{K}^\bullet \otimes_{\mathcal{O}_X} \mathcal{N}^\bullet)
\subset
\text{Tot}(\mathcal{M}^\bullet \otimes_{\mathcal{O}_X} \mathcal{N}^\bullet)$
by flatness of the modules $\mathcal{N}^n$, we see that $\eta$ factors through
$$
\tilde \eta :
\mathcal{O}_X \to
\text{Tot}(\mathcal{K}^\bullet \otimes_{\mathcal{O}_X} \mathcal{N}^\bullet)
$$
Denoting $K$ the object of $D(\mathcal{O}_X)$ represented by
$\mathcal{K}^\bullet$ we find a commutative diagram
$$
\xymatrix{
M \ar[rr]_-{\eta \otimes 1} \ar[rrd]_{\tilde \eta \otimes 1} & &
M \otimes^\mathbf{L} N \otimes^\mathbf{L} M
\ar[r]_-{1 \otimes \epsilon} &
M \\
& &
K \otimes^\mathbf{L} N \otimes^\mathbf{L} M
\ar[u] \ar[r]^-{1 \otimes \epsilon} &
K \ar[u]
}
$$
Since the composition of the upper row is the identity on $M$
we conclude that $M$ is a direct summand of $K$ in $D(\mathcal{O}_X)$.

\medskip\noindent
As a first use of the argument above, we can choose the subcomplex
$\mathcal{K}^\bullet = \sigma_{\geq a} \tau_{\leq b}\mathcal{M}^\bullet$
with $a < n_i < b$ for $i = 1, \ldots, N$. Thus $M$ is a direct
summand in $D(\mathcal{O}_X)$ of a bounded complex and we conclude
we may assume $M$ is in $D^b(\mathcal{O}_X)$. (Recall that the process
above involves shrinking $X$.)

\medskip\noindent
Since $M$ is in $D^b(\mathcal{O}_X)$ we may choose
$\mathcal{M}^\bullet$ to be a bounded above complex of
flat modules (by Modules, Lemma \ref{modules-lemma-module-quotient-flat} and
Derived Categories, Lemma \ref{derived-lemma-subcategory-left-resolution}).
Then we can choose $\mathcal{K}^\bullet = \sigma_{\geq a}\mathcal{M}^\bullet$
with $a < n_i$ for $i = 1, \ldots, N$ in the argument above.
Thus we find that we may assume $M$ is a direct summand in
$D(\mathcal{O}_X)$ of a bounded complex of flat modules.
In particular,  $M$ has finite tor amplitude.

\medskip\noindent
Say $M$ has tor amplitude in $[a, b]$. Assuming $M$ is $m$-pseudo-coherent
we are going to show that (after shrinking $X$) we may assume $M$
is $(m - 1)$-pseudo-coherent. This will finish the proof by
Lemma \ref{lemma-perfect-precise} and the fact that
$M$ is $(b + 1)$-pseudo-coherent in any case.
After shrinking $X$ we may assume there exists a strictly perfect
complex $\mathcal{E}^\bullet$ and a map $\alpha : \mathcal{E}^\bullet \to M$
in $D(\mathcal{O}_X)$ such that $H^i(\alpha)$ is an isomorphism for
$i > m$ and surjective for $i = m$. We may and do assume
that $\mathcal{E}^i = 0$ for $i < m$. Choose a distinguished triangle
$$
\mathcal{E}^\bullet \to M \to L \to \mathcal{E}^\bullet[1]
$$
Observe that $H^i(L) = 0$ for $i \geq m$. Thus we may represent
$L$ by a complex $\mathcal{L}^\bullet$ with $\mathcal{L}^i = 0$
for $i \geq m$. The map $L \to \mathcal{E}^\bullet[1]$
is given by a map of complexes
$\mathcal{L}^\bullet \to \mathcal{E}^\bullet[1]$
which is zero in all degrees except in degree $m - 1$
where we obtain a map $\mathcal{L}^{m - 1} \to \mathcal{E}^m$, see
Derived Categories, Lemma \ref{derived-lemma-negative-exts}.
Then $M$ is represented by the complex
$$
\mathcal{M}^\bullet :
\ldots \to
\mathcal{L}^{m - 2} \to
\mathcal{L}^{m - 1} \to
\mathcal{E}^m \to
\mathcal{E}^{m + 1} \to \ldots
$$
Apply the discussion in the second paragraph to this complex to get
sections $f_i$ of $\mathcal{M}^{n_i}$ for $i = 1, \ldots, N$.
For $n < m$ let $\mathcal{K}^n \subset \mathcal{L}^n$
be the $\mathcal{O}_X$-submodule generated by the sections
$f_i$ for $n_i = n$ and $d(f_i)$ for $n_i = n - 1$.
For $n \geq m$ set $\mathcal{K}^n = \mathcal{E}^n$.
Clearly, we have a morphism of
distinguished triangles
$$
\xymatrix{
\mathcal{E}^\bullet \ar[r] &
\mathcal{M}^\bullet \ar[r] &
\mathcal{L}^\bullet \ar[r] &
\mathcal{E}^\bullet[1] \\
\mathcal{E}^\bullet \ar[r] \ar[u] &
\mathcal{K}^\bullet \ar[r] \ar[u] &
\sigma_{\leq m - 1}\mathcal{K}^\bullet \ar[r] \ar[u] &
\mathcal{E}^\bullet[1] \ar[u]
}
$$
where all the morphisms are as indicated above.
Denote $K$ the object of $D(\mathcal{O}_X)$ corresponding to the complex
$\mathcal{K}^\bullet$.
By the arguments in the second paragraph of the proof we obtain
a morphism $s : M \to K$ in $D(\mathcal{O}_X)$ such that the composition
$M \to K \to M$ is the identity on $M$. We don't know that the
diagram
$$
\xymatrix{
\mathcal{E}^\bullet \ar[r] &
\mathcal{K}^\bullet \ar@{=}[r] &
K \\
\mathcal{E}^\bullet \ar[u]^{\text{id}} \ar[r]^i &
\mathcal{M}^\bullet \ar@{=}[r] &
M \ar[u]_s
}
$$
commutes, but we do know it commutes after composing with the
map $K \to M$. By Lemma \ref{lemma-local-actual} after shrinking $X$ we may
assume that $s \circ i$ is given by a map of complexes
$\sigma : \mathcal{E}^\bullet \to \mathcal{K}^\bullet$.
By the same lemma we may assume the composition of $\sigma$
with the inclusion $\mathcal{K}^\bullet \subset \mathcal{M}^\bullet$
is homotopic to zero by some homotopy
$\{h^i : \mathcal{E}^i \to \mathcal{M}^{i - 1}\}$.
Thus, after replacing $\mathcal{K}^{m - 1}$ by
$\mathcal{K}^{m - 1} + \Im(h^m)$ (note that after doing this
it is still the case that $\mathcal{K}^{m - 1}$ is generated
by finitely many global sections), we see that
$\sigma$ itself is homotopic to zero!
This means that we have a commutative solid diagram
$$
\xymatrix{
\mathcal{E}^\bullet \ar[r] &
M \ar[r] &
\mathcal{L}^\bullet \ar[r] &
\mathcal{E}^\bullet[1] \\
\mathcal{E}^\bullet \ar[r] \ar[u] &
K \ar[r] \ar[u] &
\sigma_{\leq m - 1}\mathcal{K}^\bullet \ar[r] \ar[u] &
\mathcal{E}^\bullet[1] \ar[u] \\
\mathcal{E}^\bullet \ar[r] \ar[u] &
M \ar[r] \ar[u]^s &
\mathcal{L}^\bullet \ar[r] \ar@{..>}[u] &
\mathcal{E}^\bullet[1] \ar[u]
}
$$
By the axioms of triangulated categories we obtain a dotted
arrow fitting into the diagram.
Looking at cohomology sheaves in degree $m - 1$ we see that we obtain
$$
\xymatrix{
H^{m - 1}(M) \ar[r] &
H^{m - 1}(\mathcal{L}^\bullet) \ar[r] &
H^m(\mathcal{E}^\bullet) \\
H^{m - 1}(K) \ar[r] \ar[u] &
H^{m - 1}(\sigma_{\leq m - 1}\mathcal{K}^\bullet) \ar[r] \ar[u] &
H^m(\mathcal{E}^\bullet) \ar[u] \\
H^{m - 1}(M) \ar[r] \ar[u] &
H^{m - 1}(\mathcal{L}^\bullet) \ar[r] \ar[u] &
H^m(\mathcal{E}^\bullet) \ar[u]
}
$$
Since the vertical compositions are the identity in both the
left and right column, we conclude the vertical composition
$H^{m - 1}(\mathcal{L}^\bullet) \to
H^{m - 1}(\sigma_{\leq m - 1}\mathcal{K}^\bullet) \to
H^{m - 1}(\mathcal{L}^\bullet)$ in the middle is surjective!
In particular $H^{m - 1}(\sigma_{\leq m - 1}\mathcal{K}^\bullet) \to
H^{m - 1}(\mathcal{L}^\bullet)$ is surjective.
Using the induced map of long exact sequences of cohomology
sheaves from the morphism of triangles above, a diagram chase
shows this implies $H^i(K) \to H^i(M)$ is an isomorphism
for $i \geq m$ and surjective for $i = m - 1$.
By construction we can choose an $r \geq 0$ and a surjection
$\mathcal{O}_X^{\oplus r} \to \mathcal{K}^{m - 1}$. Then the
composition
$$
(\mathcal{O}_X^{\oplus r} \to \mathcal{E}^m \to
\mathcal{E}^{m + 1} \to \ldots ) \longrightarrow
K \longrightarrow M
$$
induces an isomorphism on cohomology sheaves in degrees $\geq m$ and
a surjection in degree $m - 1$ and the proof is complete.
\end{proof}

\end{document}
